Un élève de \qty{\v{m}}{kg} monte en courant les escaliers de l'école, d'une hauteur totale de \qty{\v{h}}{m}, en \qty{\v{t}}{s}. On prend $g = \qty{9.81}{m/s^2}$.

\begin{enumerate}[label=\alph*)]
    \item Quelle énergie potentielle gravitationnelle l'élève gagne-t-il ?
    \item Quelle puissance moyenne développe-t-il pour monter ?
    \item Pendant combien de temps cette énergie permettrait-elle d'allumer une ampoule de \qty{\v{P}}{W} ?
\end{enumerate}
